\section{The unified core}\label{sec:core}

\subsection{Block-head candidate parameterisation}

Let
\begin{equation*}
  r=\frac{A_j+5^j q}{2^{\wordWeight_j}}
\end{equation*}
be a block-head state.  The unified core writes the predecessor
candidate in the form
\begin{equation*}
  x=2^{e-1}g+\delta\,5^{j-1},
\qquad g=\Orbit_7(j-1),
\end{equation*}
where $e\ge2$ is the incoming step weight, $\delta\in\{1,3\}$ is the
reset parameter, and the reset step has weight $t\in\{1,2\}$.  The exact
relations are:
\begin{enumerate}[leftmargin=*]
\item the reset equation
  $\ResetWindow(j,k_0,t,\delta,s)$ fixes $s,j,k_0,t,\delta$;
\item the first-block terminal equation gives
  $s-1=2^{e-1}g$;
\item the candidate satisfies
  $x=2^{e-1}g+\delta 5^{j-1}$.
\end{enumerate}

These statements are formalised in the Lean bridge module
\texttt{UnifiedCoreBridge.lean}: the reset-predecessor lemma, the
candidate parameterisation, and the first-block terminal identity.

\subsection{Segment-length exclusions}

\begin{lemma}[Segment-length exclusions]\label{lem:segments}
Let $d$ be the length of the full-orbit segment from $g$ to the
candidate $x$.  Then $d\ne1$, $d\ne2$, $d\ne3$, and $d\ge4$ is
impossible.  Hence the candidate family is empty.
\end{lemma}

\begin{remark}[Status of \Cref{lem:segments}]
The mathematical proof is closed.  The Lean side has compiled the
$d=1$, $d=2$, $d=3$, and $d\ge4$ exclusions, together with the
$e=3$, $j=17$ modular exclusions.
\end{remark}

\begin{center}
\begin{tabular}{lll}
\toprule
Statement & Math & Lean \\
\midrule
Candidate parameterisation & proven & compiled \\
Full-orbit to candidate derivation & proven & pending \\
Reset-head predecessor & proven & compiled \\
First-block terminal equation & proven & compiled \\
$d=1$ exclusion & proven & compiled \\
$d=2$ exclusion & proven & compiled \\
$d=3$ family exclusion & proven & compiled \\
$d\ge4$ exclusion & proven & compiled \\
Unified-core final inequality & proven & pending \\
\bottomrule
\end{tabular}
\end{center}

\subsection{Word rigidity and the first block}

The full-orbit reachability of the block head forces the first block to
have the rigid shape
\begin{equation*}
  y\longrightarrow x\longrightarrow r_j\longrightarrow z
  \longrightarrow w,
\end{equation*}
with $k_0=0$ and the first-block terminal equation
\begin{equation*}
  s-1=2^{e-1}g.
\end{equation*}
The Lean side has compiled the first-block terminal identity and the
candidate parameterisation.  The complete derivation from the
block-head premises and $\FullOrbit(r)$ to this candidate
parameterisation is mathematically closed but is still marked as Lean
pending.

\subsection{The q0 interval and block-head identity}

The block-head numerator satisfies the exact identity
\begin{equation*}
  A_j+5^jq=2^L\left(B+\delta5^j\right),
\qquad 0<B<5^j,\quad A_j<5^j,
\end{equation*}
which forces
\begin{equation*}
  q=m+\delta\,2^L,\qquad m<2^L.
\end{equation*}
This is the explicit $q_0$ form.  Under the reset equation, the same
numerator has the block-head form
\begin{equation*}
  A_j+5^jq
    =2^{W_p}\left(5^{k+1}s-4+\delta5^j\right),
\end{equation*}
which is the block-head identity of a reset.  The interval
condition
\begin{equation*}
  2^{W_p}\le q<2^{W_j}
\end{equation*}
is equivalent to the block-head bounds
\begin{equation*}
  5^j\le2^t r_j,\qquad r_j<5^j,
\end{equation*}
by the $q_0$-interval criterion.  These three statements are
compiled.

\subsection{Segment equations}

Let $u_1,\dots,u_d$ be the step weights of the full-orbit segment from
$g=\Orbit_7(j-1)$ to $x=\Orbit_7(j-1+d)$, and let
$W=u_1+\cdots+u_d$.  The exact segment equation is
\begin{equation*}
  2^{W+e-1}g+2^W\delta5^{j-1}
    =5^d\,g+\sum_{i=0}^{d-1}2^{W_i}5^{d-1-i},
\end{equation*}
where $W_i$ is the prefix weight inside the segment, $W=W_d$, and $e$
is the incoming step weight.  This is the candidate form of the Lean
equation
\begin{equation*}
  2^W\,x=5^d\,g+\wordA(w)
\end{equation*}
obtained by substituting $x=2^{e-1}g+\delta5^{j-1}$.

The derivation is direct: the segment word $w$ has length $d$ and maps
$g$ to $x$, so the exact word-orbit identity gives
$2^Wx=5^dg+\wordA(w)$.  Substitution of the candidate parameterisation
is exactly the segment-candidate equation formalised in Lean.  In the
special cases
$d=1$ and $d=2$, the word molecules are $\wordA=1$ and
$\wordA=5+2^{u_1}$.

For $d=1$ the equation reduces to
\begin{equation*}
  5g+1=2^{1+4a}x,
\end{equation*}
which is impossible by the size bound $g<5^{j-1}/4$.  For $d=2$ the
surviving branch forces the two modular equations
\begin{equation*}
  5^{j-1}\equiv6\pmod{17},\qquad
  5^{j-1}\equiv45\pmod{256},
\end{equation*}
which contradict each other modulo $16$.  For $d=3$ the only surviving
family has first big step at depth $j+4$ with weight $6$, contradicting
the finite base.  For $d\ge4$, the first big step is either at depth
$j-2$ with weight $3$, or later with weight at least $5$, again
contradicting the finite base.

The orbit-level $d=1$ and $d=2$ exclusions are compiled in
\texttt{UnifiedCoreBridge.lean}.

\subsection{Detailed exclusion arguments}

\paragraph{$d=1$.}
The segment equation is
\begin{equation*}
  5g+1=2^{1+4a}\,(2^{e-1}g+\delta5^{j-1}).
\end{equation*}
For $e\ge3$ or $a\ge1$, the coefficient $2^{e+4a}$ is at least $8$,
while the right-hand side contains a positive multiple of
$2^{1+4a}\delta5^{j-1}$, forcing a contradiction with the size bound
$g<5^{j-1}/4$.  The only remaining case is $e=2,a=0$, which gives
\begin{equation*}
  g+1=2\delta\,5^{j-1},
\end{equation*}
so $g\ge2\cdot5^{j-1}-1$ for $\delta\ge1$, again contradicting
$g<5^{j-1}/4$.

\paragraph{$d=2$.}
The size constraint $2^{1+4a}\delta\le24$ forces $a=0$.  The segment
equation is
\begin{equation*}
  (2^{u_1+e}-25)g+2^{u_1+1}\delta5^{j-1}=5+2^{u_1}.
\end{equation*}
The only surviving branch is $\delta=1$, $e=2$, $u_1=1$; substituting
into the segment equation gives
\begin{equation*}
  (8-25)g+4\cdot5^{j-1}=7,
\end{equation*}
hence
\begin{equation*}
  17g=4\cdot5^{j-1}-7,
\end{equation*}
hence $j-1\equiv3\pmod{16}$.  The candidate congruence
$x\equiv743\pmod{1280}$ then forces
\begin{equation*}
  5^{j-1}\equiv45\pmod{256},
\end{equation*}
so $j-1\equiv7\pmod{64}$, contradicting $j-1\equiv3\pmod{16}$.

\subsection{The $d=2$ survivor in detail}

The survivor equations are
\begin{equation*}
  17g=4\cdot5^{j-1}-7,\qquad
  x\equiv743\pmod{1280},\qquad x=2g+5^{j-1}.
\end{equation*}
Modulo $17$, the first equation gives $4\cdot5^{j-1}\equiv7$, so
\begin{equation*}
  5^{j-1}\equiv6\pmod{17}.
\end{equation*}
Since $5^3\equiv6\pmod{17}$ and the order of $5$ modulo $17$ is $16$,
we get $j-1\equiv3\pmod{16}$.

Eliminating $g$ gives
\begin{equation*}
  17x=25\cdot5^{j-1}-14,
\end{equation*}
so $x\equiv743\pmod{1280}$ forces
\begin{equation*}
  25\cdot5^{j-1}\equiv1125\pmod{1280}.
\end{equation*}
Dividing by $5$ and multiplying by the inverse of $5$ modulo $256$
reduces this to
\begin{equation*}
  5^{j-1}\equiv45\pmod{256}.
\end{equation*}
Since $5^7\equiv45\pmod{256}$ and the order of $5$ modulo $256$ is
$64$, we get $j-1\equiv7\pmod{64}$.  The two congruences contradict
each other modulo $16$.

\paragraph{$d=3$.}
The remaining family is
\begin{equation*}
  t_j=2,\quad \delta=1,\quad e=2,\quad u_1=u_2=1,\quad a=0,
\end{equation*}
so $u_3=1$; the segment equation becomes
\begin{equation*}
  16g+8\cdot5^{j-1}=125g+39,
\end{equation*}
which gives
\begin{equation*}
  g=\frac{8\cdot5^{j-1}-39}{109},\qquad
  j-1\equiv923\pmod{1728}.
\end{equation*}
The word shape is
\begin{equation*}
  g\to x\to r_j\to z\to w,
\end{equation*}
and the first step of weight at least $3$ is the final step at depth
$j+4$ with weight $6$.  The finite base gives first big weight $3$ at
depth $15$, a contradiction.

\paragraph{$d\ge4$.}
The first big step is one of three types.  If $e\ge3$, the first big
step is $g_{\mathrm{prev}}\to g$; the finite base forces $e=3$ and
$j=17$, and the candidate residue modulo $640$ or $1280$ is
excluded.  If $e=2$ and $a\ge1$, the first big step is $y\to x$ with
weight $1+4a\ge5$.  If $e=2$ and $a=0$, the first big step is
$z\to w$ with weight $5$ or $6$.  All three contradict the finite base.

\subsection{Branch ledger}

\begin{center}
\small
\begin{tabular}{@{}p{0.11\textwidth}p{0.42\textwidth}p{0.38\textwidth}@{}}
\toprule
Case & Surviving branch & Contradiction \\
\midrule
$d=1$ & $e=2$, $a=0$ &
  $g=2\cdot5^{j-1}-1$ vs.\ $g<5^{j-1}/4$ \\
$d=2$ & $\delta=1$, $e=2$, $u_1=1$ &
  $j-1\equiv3\pmod{16}$ and $j-1\equiv7\pmod{64}$ \\
$d=3$ & $t_j=2$, $\delta=1$, $e=2$, $u_1=u_2=1$, $a=0$ &
  first big step at depth $j+4$ has weight $6$, not $3$ \\[2pt]
$d\ge4$ & three big-step types &
  weight $\ge5$, or residue exclusion for $e=3$, $j=17$ \\
\bottomrule
\end{tabular}
\end{center}

Each row records the exact contradiction used in the mathematical
proof.  The Lean names for the individual exclusions are listed in
the status tables above.

\subsection{Modular arithmetic ledger}

\begin{center}
\small
\begin{tabular}{@{}p{0.10\textwidth}p{0.56\textwidth}p{0.28\textwidth}@{}}
\toprule
Case & Congruence and role & Lean \\
\midrule
$d=1$ & no congruence; the size bound $g<5^{j-1}/4$ suffices &
  \texttt{d1\_\allowbreak exclusion} \\
$d=2$ & $5^{j-1}\equiv6\pmod{17}$ and $x\equiv743\pmod{1280}$ force
  $j-1\equiv3\pmod{16}$ and $5^{j-1}\equiv45\pmod{256}$ &
  \texttt{d2\_\allowbreak exclusion} \\
$d=3$ & $j-1\equiv923\pmod{1728}$ fixes the unique surviving family &
  \texttt{d3\_\allowbreak family\_\allowbreak big\_\allowbreak
  weight\_\allowbreak excluded} \\
$d\ge4$ & the $e=3$, $j=17$ branch is excluded by residues modulo
  $640$ or $1280$ &
  \texttt{dge4\_\allowbreak e3\_\allowbreak j17\_\allowbreak t1\_\allowbreak
  excluded}, \texttt{dge4\_\allowbreak e3\_\allowbreak j17\_\allowbreak
  t2\_\allowbreak delta3\_\allowbreak excluded} \\
\bottomrule
\end{tabular}
\end{center}

These congruences are exact and are the only modular inputs needed
after the size estimates; they are proved inside the corresponding
segment exclusions.

\subsection{Worked modular ledger}

The three congruence families used by the exclusions are:
\begin{align*}
  5^{j-1}&\equiv6\pmod{17}
  &&\Longrightarrow j-1\equiv3\pmod{16},\\
  5^{j-1}&\equiv45\pmod{256}
  &&\Longrightarrow j-1\equiv7\pmod{64},\\
  5^{j-1}&\equiv73\pmod{109}
  &&\text{for the }d=3\text{ family}.
\end{align*}
The first two follow from $5^3\equiv6\pmod{17}$ and
$5^7\equiv45\pmod{256}$ with the respective orders $16$ and $64$.
The third is the integrality condition for
\begin{equation*}
  g=\frac{8\cdot5^{j-1}-39}{109}.
\end{equation*}
Its consistency is checked by repeated squaring modulo $109$:
\begin{center}
\small
\begin{tabular}{@{}rrrr@{}}
\toprule
$k$ & $5^k\bmod 109$ & $k$ & $5^k\bmod 109$ \\
\midrule
1 & 5 & 64 & 97 \\
2 & 25 & 128 & 35 \\
4 & 80 & 256 & 26 \\
8 & 78 & 512 & 22 \\
16 & 89 & 923 & 73 \\
32 & 73 & & \\
\bottomrule
\end{tabular}
\end{center}
The last entry uses
\begin{equation*}
  923=512+256+128+16+8+2+1
\end{equation*}
and the corresponding product of residues, giving
\begin{equation*}
  5^{923}\equiv73\pmod{109},
\qquad
  8\cdot73=584\equiv39\pmod{109}.
\end{equation*}
These checks are part of the compiled exclusions; the paper records
them here so that the modular layer can be read without a calculator.

\subsection{Block-head candidate exclusion}

The segment-length exclusions show that no candidate
\begin{equation*}
  x=2^{e-1}g+\delta5^{j-1}
\end{equation*}
can satisfy the full set of constraints coming from
\texttt{All36\_20PremisesNoHge} and $\FullOrbit(r)$.  Hence no block
head satisfies the candidate family of document 36.30.7.  This is the
mathematical closure of the unified core's candidate layer.  The Lean
side has compiled the generic exclusion theorems, while the complete
derivation from the no-$H_{\ge}$ premises to those exclusions remains
pending.

\subsection{Proof outline}

\begin{enumerate}[leftmargin=*]
\item Extract the block head
  $r=(A_j+5^jq)/2^{\wordWeight_j}$ from the no-$H_{\ge}$ premises.
\item Use $\FullOrbit(r)$ to place the block head on the exact orbit of
  $7$ and obtain the incoming weight $e$ and the state
  $g=\Orbit_7(j-1)$.
\item Parameterise the predecessor by
  $x=2^{e-1}g+\delta5^{j-1}$.
\item Write the exact segment equation for $d$ steps from $g$ to $x$.
\item Exclude $d=1$, $d=2$, $d=3$, and $d\ge4$; hence the candidate
  family is empty.
\item Pass through the CRT candidate $r_{j,0}$ and the residue $y^*$ to
  obtain the final valuation inequality.
\end{enumerate}

\subsection{From candidate emptiness to the final inequality}

The exclusion of every candidate block head leaves the CRT candidate
$r_{j,0}$ and the terminal residue $y^*$.  The terminal inequality
\begin{equation*}
  \vtwo{5^{L+3}w_{\mathrm{term}}+1}\le H_s-2
\end{equation*}
is equivalent to $y^*\ne0$, and the remaining Lean statement
\texttt{unified\_core\_final\_no\_hge} is exactly this valuation inequality
under the no-$H_{\ge}$ premises and the full-orbit reachability of the
block head.  The mathematical layer is closed; the Lean statement is
pending.

\subsection{The remaining Lean statement}

The exact statement of the pending theorem is:
\begin{verbatim}
theorem unified_core_final_no_hge
    (j Wp Wj q Aj A_s s W_s r_s L H_s : Nat)
    (weight : Nat -> Nat) (r : Nat)
    (hPrem : All36_20PremisesNoHge
       j Wp Wj q Aj A_s s W_s r_s L H_s weight)
    (hrj : r = (Aj + 5 ^ j * q) / 2 ^ Wj)
    (hReach : S6Audit.FullOrbitFrom7 r)
    (hReset : exists s0 k t delta r_prev : Nat,
      S6Audit.ResetHeadEq s0 j k t delta r /\
        s0 * 5 ^ k = r_prev + 1)
    (hH : 2 <= H_s) :
    twoValuation (5 ^ (L + 3) * wTerminal L r_s + 1) <= H_s - 2 := by
  sorry
\end{verbatim}

The reset premise is arithmetic, not a general-orbit premise: it
records the exact reset equation and the previous-terminal equation
$s_0\cdot5^k=r_{\mathrm{prev}}+1$.  The block head is required to lie
on the real accelerated orbit by $\FullOrbit(r)$.

The premises are: the no-$H_{\ge}$ record, the block-head equality
$r=(A_j+5^jq)/2^{W_j}$, the full-orbit reachability of $r$, and the
domain condition $2\le H_s$.  The conclusion is the terminal valuation
inequality of \Cref{thm:core}.  This is the single \texttt{sorry}
remaining in the unified-core audit file.

\subsection{The CRT candidate and the terminal residue}

Let $n=s-j$ and $\Delta=W_s-W_j$.  The block molecule $B$ of the suffix
satisfies
\begin{equation*}
  3B+2^\Delta\le 3\cdot5^n-2\cdot4^n,
\end{equation*}
and the CRT candidate $r_{j,0}$ is the least nonnegative solution of
the exact integer equation
\begin{equation*}
  3\cdot5^n\,r_{j,0}+3B+2^\Delta
    =2^{\Delta+L+4}\left(u+m'\,2^{H_s-1}\right),
\end{equation*}
where $u$ is the least nonnegative residue of
$-5^{-(L+3)}$ modulo $2^{H_s-1}$ and $m'\ge0$.

The terminal odd part is
\begin{equation*}
  w_{\mathrm{term}}=\frac{3r_s+1}{2^{L+4}},
\end{equation*}
and the residue
\begin{equation*}
  y^*=
  -\left(w_{\mathrm{term}}+5^{-(L+3)}\right)
  \left(3\cdot5^s\right)^{-1}
  \pmod{2^{H_s-1}}
\end{equation*}
is the obstruction to the final valuation inequality.  The Lean
equivalences
\begin{equation*}
  \vtwo{5^{L+3}w_{\mathrm{term}}+1}\le H_s-2
  \iff y^*\ne0
\end{equation*}
are proved in
\texttt{terminal\_bound\_iff\_yStar\_ne\_zero}.  The remaining Lean
statement \texttt{unified\_core\_final\_no\_hge} is the passage from
the no-$H_{\ge}$ premises and $\FullOrbit(r)$ to this inequality.

\subsection{Derivation of the CRT equation}

The CRT equation is not introduced as an independent congruence; it is
the failure equation of the terminal valuation written in block
coordinates.  Let $n=s-j$ and $\Delta=W_s-W_j$, and let $B$ be the
block molecule of the suffix of length $L$.  The suffix is a valid
word from $r_{j,0}$ to $r_s$, so the exact word-orbit identity gives
\begin{equation*}
  2^{\Delta}\,r_s=5^n\,r_{j,0}+B.
\end{equation*}
The terminal odd part is defined by
\begin{equation*}
  w_{\mathrm{term}}=\frac{3r_s+1}{2^{L+4}},
\end{equation*}
so, multiplying the tail equation by $3\cdot2^{\Delta}$ and adding
$2^{\Delta}$,
\begin{equation*}
  2^{\Delta+L+4}\,w_{\mathrm{term}}
    =2^{\Delta}(3r_s+1)
    =3\cdot5^n\,r_{j,0}+3B+2^{\Delta}.
\end{equation*}
This single identity is the bridge between the orbit equation and the
valuation inequality.

The valuation inequality
\begin{equation*}
  \vtwo{5^{L+3}w_{\mathrm{term}}+1}\le H_s-2
\end{equation*}
fails exactly when
\begin{equation*}
  5^{L+3}w_{\mathrm{term}}+1\equiv0\pmod{2^{H_s-1}},
\end{equation*}
i.e.\ exactly when
\begin{equation*}
  w_{\mathrm{term}}\equiv-5^{-(L+3)}
  \pmod{2^{H_s-1}}.
\end{equation*}
The inverse exists because $5$ is odd.  Let $u$ be the least
nonnegative residue of $-5^{-(L+3)}$ modulo $2^{H_s-1}$.  Substituting
the failure congruence into the identity above gives the CRT equation
\begin{equation*}
  3\cdot5^n\,r_{j,0}+3B+2^{\Delta}
    =2^{\Delta+L+4}\left(u+m'\,2^{H_s-1}\right)
\end{equation*}
for some integer $m'\ge0$.  This is the exact form used in
\texttt{rj0\_\allowbreak spec\_\allowbreak eq}.

The terminal residue $y^*$ is the normalised failure residue
\begin{equation*}
  y^*=
  -\left(w_{\mathrm{term}}+5^{-(L+3)}\right)
  \left(3\cdot5^s\right)^{-1}
  \pmod{2^{H_s-1}}.
\end{equation*}
By construction $y^*=0$ if and only if the failure congruence holds,
so the valuation inequality is equivalent to $y^*\ne0$.  The factor
$\left(3\cdot5^s\right)^{-1}$ is the normalisation inherited from the
orbit equation for the terminal state; it is invertible modulo
$2^{H_s-1}$ because both $3$ and $5^s$ are odd.  The equivalence is
\texttt{terminal\_\allowbreak bound\_\allowbreak iff\_\allowbreak
yStar\_\allowbreak ne\_\allowbreak zero}.

\subsection{Size conditions}

The sufficient direction uses two inverse-size conditions.  The first
is the block-tail bound
\begin{equation*}
  3B+2^\Delta\le 3\cdot5^n-2\cdot4^n,
\end{equation*}
and the second is the base-size condition
\begin{equation*}
  3\cdot5^s+\left(3\cdot5^n-2\cdot4^n\right)
    \le 2^{\Delta+L+4}\,u,
\end{equation*}
which together imply $r_{j,0}\ge5^j$.  This is proved in Lean by
\texttt{rj0\_ge\_of\_size\_conditions\_no\_hge}.  The full equivalence
through the CRT lift remains part of the unified core; it is not
claimed as an independent step.

\subsection{Why the size conditions imply $r_{j,0}\ge5^j$}

The two size conditions are exactly what the CRT equation needs in
order to lift the candidate above $5^j$.  The first condition
\begin{equation*}
  3B+2^{\Delta}\le3\cdot5^n-2\cdot4^n
\end{equation*}
controls the correction term, and the second condition
\begin{equation*}
  3\cdot5^s+\left(3\cdot5^n-2\cdot4^n\right)
    \le2^{\Delta+L+4}\,u
\end{equation*}
controls the leading term.  From the CRT equation,
\begin{equation*}
  3\cdot5^n\,r_{j,0}
    =2^{\Delta+L+4}\left(u+m'2^{H_s-1}\right)-3B-2^{\Delta}.
\end{equation*}
The right-hand side is at least
\begin{equation*}
  2^{\Delta+L+4}\,u-\left(3\cdot5^n-2\cdot4^n\right),
\end{equation*}
and the second condition bounds this below by $3\cdot5^s$.  Therefore
\begin{equation*}
  r_{j,0}\ge\frac{3\cdot5^s}{3\cdot5^n}=5^{s-n}=5^j,
\end{equation*}
because $n=s-j$.  This is the arithmetic content of
\texttt{rj0\_\allowbreak ge\_\allowbreak of\_\allowbreak
size\_\allowbreak conditions\_\allowbreak no\_\allowbreak hge}; the Lean
statement packages the divisibility and nonnegativity facts that the
paper suppresses.

The lower bound $r_{j,0}\ge5^j$ is then compared with the interval
constraints carried by the no-$H_{\ge}$ premises.  The passage from
the block-head interval to the valuation inequality is exactly the
remaining Lean bridge \texttt{unified\_\allowbreak core\_\allowbreak
final\_\allowbreak no\_\allowbreak hge}; the paper does not claim it as
closed.

\subsection{Relation to source documents}

\begin{center}
\small
\begin{tabular}{@{}p{0.20\textwidth}p{0.36\textwidth}p{0.38\textwidth}@{}}
\toprule
Document & Paper & Lean \\
\midrule
36.30.22 & candidate parameterisation & \texttt{UnifiedCoreBridge} compiled \\
36.30.23 & segment exclusions & \texttt{FinitePrefix} and \\
         &                      & \texttt{UnifiedCoreBridge}; compiled \\
36.30.7 & block-head candidate exclusion & pending full derivation \\
36.20 & final valuation inequality & \texttt{UnifiedCoreAudit} pending \\
\bottomrule
\end{tabular}
\end{center}

\subsection{The finite base}

The full orbit is explicitly expanded for the first $17$ states:
\begin{equation*}
  7,9,23,29,73,183,229,573,1433,3583,4479,5599,6999,8749,21873,54683,34177.
\end{equation*}
The first step of weight $\ge3$ is the step from depth $15$ to depth
$16$, and its weight is exactly $3$.  This is proved in
\texttt{FinitePrefix.lean} by explicit rewriting, with no
\texttt{native\_decide}.

The expansion stops at 17 states because every exclusion either uses
the first big step (depth $15$, weight $3$) or reduces the branch to
$j=17$.  No later full-orbit data is needed.

\subsection{How the finite base is used}

\begin{center}
\small
\begin{tabular}{@{}p{0.42\textwidth}p{0.32\textwidth}p{0.20\textwidth}@{}}
\toprule
Finite fact & Lean name & Used in \\
\midrule
first 17 states expanded &
  \texttt{fullOrbitIter\_\allowbreak prefix\_expand} & all exclusions \\
step weights for $n<16$ &
  \texttt{fullOrbit\_\allowbreak prefix\_step\_weights} & all exclusions \\
first $t\ge3$ is at $n=15$ and equals $3$ &
  \texttt{fullOrbit\_\allowbreak first\_t\_ge3\_is\_exactly\_3} & $d=3$, $d\ge4$ \\
first big step is unique in the prefix &
  \texttt{first\_\allowbreak big\_step\_unique} & $d=3$, $d\ge4$ \\
candidate first big weight $\ne3$ &
  \texttt{candidate\_\allowbreak first\_\allowbreak big\_\allowbreak
  step\_\allowbreak weight\_\allowbreak ne\_three} & $d=3$ \\
candidate first big weight $\ge5$ &
  \texttt{candidate\_\allowbreak first\_\allowbreak big\_\allowbreak
  step\_\allowbreak weight\_\allowbreak ge\_five} & $d\ge4$ \\
\bottomrule
\end{tabular}
\end{center}

The finite base is used only to identify the first big step and its
weight.  Every later step is exact arithmetic.  No probabilistic data
or \texttt{native\_decide} enters the chain.

\subsection{Final valuation inequality}

\begin{theorem}[Unified core, mathematical form]\label{thm:core}
Under the no-$H_{\ge}$ premises of document 36.20, the block-head
reachability $\FullOrbit(r)$, and the domain condition $2\le H_s$, we
have
\begin{equation*}
  \vtwo{5^{L+3}\,w_{\mathrm{term}}+1}\le H_s-2,
\end{equation*}
where $w_{\mathrm{term}}=(3r_s+1)/2^{L+4}$.
\end{theorem}

\begin{remark}[Formalisation status]
\statusmath{}
\statuspending{} for
\texttt{UnifiedCoreAudit.unified\_core\_final\_no\_hge}.
The Lean statement is the unique \texttt{sorry} in the unified-core
audit file.
The paper does not claim Lean closure of \Cref{thm:core}.
\end{remark}

\subsection{Corrected decisive window bound}

The old arithmetic window bound
\begin{equation*}
  \vtwo{5^{k_0+1}s+\delta 5^j}\le 2j+10
\end{equation*}
is false without reachability constraints: there is an arithmetic
counterexample
\begin{equation*}
  (j,k_0,t,\delta,s)=(36,0,2,3,\,2^{83}-3\cdot5^{35})
\end{equation*}
with valuation $83>82$.  The corrected bound therefore carries the
predicate $\ResetWindow(j,k_0,t,\delta,s)$, which fixes the same
$j,k_0,t$ in the reset equation, the size bound, the previous-terminal
reachability, and the full-orbit reachability of the reset head.

\subsection{Proof-state ledger}

The following matrix records the status of every layer of the main
chain.  The columns are the mathematical status and the Lean status;
the rows follow the dependency order of the proof.

\begin{center}
\small
\begin{tabular}{@{}p{0.38\textwidth}p{0.18\textwidth}p{0.34\textwidth}@{}}
\toprule
Layer & Math & Lean \\
\midrule
word orbit identity & proven & compiled \\
PMI identity and PMI-B budgets & proven & compiled \\
finite-prefix base & proven & compiled \\
candidate parameterisation & proven & compiled \\
reset-head predecessor & proven & compiled \\
first-block terminal equation & proven & compiled \\
$d=1$ exclusion & proven & compiled \\
$d=2$ exclusion & proven & compiled \\
$d=3$ family exclusion & proven & compiled \\
$d\ge4$ exclusion & proven & compiled \\
CRT equation and candidate & proven & compiled \\
terminal residue equivalence & proven & compiled \\
size-condition lower bound & proven & compiled \\
final valuation inequality & proven & pending \\
cycle-word extraction and QB-8 split & proven & compiled \\
corrected window bound & proven & compiled (conditional) \\
failure-window contradiction & proven & compiled \\
no-cycle bridge & proven & compiled (conditional) \\
final assembly & proven & pending \\
\bottomrule
\end{tabular}
\end{center}

Two entries are pending: the final valuation inequality
\texttt{unified\_\allowbreak core\_\allowbreak final\_\allowbreak
no\_\allowbreak hge} and the final assembly
\texttt{five\_\allowbreak x\_\allowbreak plus\_\allowbreak one\_\allowbreak
diverges\_\allowbreak at\_\allowbreak 7}.  Every other row has a
compiled Lean representative, with the conditional rows explicitly
consuming the open inputs.  The matrix is kept in sync with the
repository audit; it is not a substitute for the build.
